% Project-specific printable Blank companion to the enterprise discovery technical reference.
% Scope: organization-wide GitHub Actions architecture discovery.
% A working form is not evidence of an approved or completed architecture.
\documentclass[11pt,letterpaper]{article}
\usepackage[T1]{fontenc}
\usepackage[utf8]{inputenc}
\usepackage{lmodern}
\usepackage[margin=0.82in,headheight=17pt]{geometry}
\usepackage{microtype}
\usepackage{booktabs,tabularx,array}
\usepackage{enumitem}
\usepackage{xcolor}
\usepackage{amssymb}
\usepackage{fancyhdr}
\usepackage{hyperref}
\definecolor{navy}{HTML}{17314C}
\definecolor{teal}{HTML}{087F78}
\hypersetup{colorlinks=true,linkcolor=navy,urlcolor=teal,pdftitle={GitHub Actions Discovery Architecture Description — Blank}}
\setlength{\parindent}{0pt}
\setlength{\parskip}{5pt}
\setlength{\emergencystretch}{3em}
\renewcommand{\arraystretch}{1.22}
\setlist[itemize]{leftmargin=*,itemsep=2pt}
\pagestyle{fancy}
\fancyhf{}
\fancyhead[L]{\small Enterprise architecture discovery}
\fancyhead[R]{\small Blank}
\fancyfoot[C]{\thepage}
\renewcommand{\headrulewidth}{0.4pt}
\newcommand{\Field}[2]{\par\noindent\textbf{#1}\par\vspace{1em}\noindent\rule{0.96\linewidth}{0.4pt}\par\vspace{1.0em}\noindent\rule{0.96\linewidth}{0.4pt}\par\vspace{1em}}
\newcommand{\Cell}[1]{\rule{0pt}{2.6em}\rule{0.86\linewidth}{0.4pt}}
\begin{document}
\begin{center}
{\LARGE\bfseries GitHub Actions Discovery Architecture Description}\par
\vspace{0.3em}
{\large Printable Blank}\par
\vspace{0.5em}
{\small GitHub Actions Architecture Discovery at Enterprise Scale}
\end{center}
\noindent\textbf{Purpose:} Describe the proposed and implemented discovery platform, its stakeholders, architectural boundaries, evidence model, views and decisions without claiming that static workflow observations prove runtime dependencies.

\noindent\textbf{Use:} Fill in the empty response areas. No example answers, review outcomes, or claims of approval are prefilled.

\noindent\textbf{Source reference:} GitHub Actions Architecture Discovery at Enterprise Scale, technical reference, October 2026.

\tableofcontents
\clearpage
\section{Identity, scope and baselines}
Establish the entity of interest, authorized organizations and intended architectural description baseline.

\Field{Entity of interest and ownership}{State platform scope, responsible service, and accountable owners.}
\Field{Authorized repository scope}{Record organizations, installations, access exclusions, archived/private coverage, and denominator.}
\Field{Project lifecycle and baseline}{State planned, pilot, deployed or transition baseline with version identifiers.}
\Field{Business purpose and exclusions}{List inventory goals and what requires external verification.}

\section{Stakeholders, concerns and requirements}
\Field{Stakeholder register}{Identify platform engineering, enterprise architects, AppSec, repository owners, privacy and operations.}
\Field{Concern and requirement IDs}{Record privacy, coverage, accuracy, provenance, latency, scale, and access concerns with IDs.}
\Field{Acceptance and verification evidence}{Link to acceptance tests, sample inventories, and requirements traceability.}
\subsection{Concern to view traceability}
\begin{tabularx}{\linewidth}{@{}p{0.31\linewidth}X@{}}
\toprule
\textbf{Record / criterion} & \textbf{Response / evidence} \\
\midrule
Concern ID and stakeholder & \Cell{Map each concern to a stakeholder and intended view.} \\
View ID, evidence and disposition & \Cell{Reference architectural view, validation evidence, and open questions.} \\
\bottomrule
\end{tabularx}


\section{System context and security boundary}
\Field{External systems and trust zones}{Describe GitHub APIs, installations, optional registries, service catalog and deployment evidence sources.}
\Field{Credential and access model}{Record least-privilege GitHub App authorization, token scope and permissions.}
\Field{Collection restrictions}{Document prohibition on untrusted action/script execution, write actions and secrets collection.}

\section{Logical processing architecture}
Discover, pin, inspect, parse, correlate and publish are distinct stages.

\Field{Collectors and API requests}{Define paginated inventory, immutable SHA pinning, backoff, and checkpointing.}
\Field{Parsers and safety constraints}{Specify GitHub Actions YAML, manifest parsing, workspace resolution, and nonexecution policy.}
\Field{Correlation and publication}{Describe typed edge creation, evidence classes, uncertainty handling and output interfaces.}
\Field{Failure paths and partial results}{Record denied, missing, truncated, malformed, stale and rate-limited states.}

\section{Architecture viewpoints and views}
\Field{Viewpoint definitions}{Describe repository topology, workflow graph, component, artifact, and deployment evidence views.}
\Field{Views and model-kind conventions}{Identify diagrams, schema versions, legends, viewpoints and correspondence rules.}
\Field{Cross-view consistency}{Describe repository-to-component, job-to-artifact and artifact-to-environment reconciliation.}

\section{Data entities and provenance}
\Field{Entity identifiers}{Define Organization, Repository, Revision, Workflow, Job, Action, Component, Artifact, Environment, Deployment.}
\Field{Evidence lineage fields}{Capture source URI/path, pinned SHA, parser and rule version, collection time, status and quality.}
\Field{E1 E2 E3 evidence semantics}{Keep observed configuration, inferred relationship, and externally verified evidence distinct.}
\Field{Storage and retention}{Record snapshot versioning, partitioning, tenancy, retention, encryption and access.}

\section{Operational deployment and quality attributes}
\Field{GitHub Action execution model}{Specify workflow triggers, runner permissions, installation authorization, input parameters and outputs.}
\Field{Capacity, bounded concurrency and recovery}{Capture request budget, cache strategy, rate-limit behavior and checkpoint recovery.}
\Field{Quality measures}{Define determinism, lineage, coverage, precision and recall with denominators.}
\Field{Publishing and access control}{Describe architecture catalog delivery, audience boundaries and data classification.}

\section{Architecture decisions and evaluation}
\Field{ADR register}{Link App access, storage, identifier, parser, confidence, cadence and retention decisions.}
\Field{Alternatives and trade-offs}{Capture selected and rejected strategies, constraints and consequences.}
\Field{Architecture evaluation and deviations}{Record quality scenarios, findings, unresolved concerns and action owners.}

\section{Traceability and change history}
\subsection{Requirements to evidence}
\begin{tabularx}{\linewidth}{@{}p{0.31\linewidth}X@{}}
\toprule
\textbf{Record / criterion} & \textbf{Response / evidence} \\
\midrule
Requirement and concern IDs & \Cell{Link objective, view, implementation, test, evidence locator and current status.} \\
Decision and view correspondence & \Cell{Connect ADRs, views, data model, deployed version and baseline.} \\
\bottomrule
\end{tabularx}

\subsection{Change control}
\begin{tabularx}{\linewidth}{@{}p{0.31\linewidth}X@{}}
\toprule
\textbf{Record / criterion} & \textbf{Response / evidence} \\
\midrule
Architecture baseline and revision & \Cell{Record change identifier, justification, approver and successor.} \\
Outstanding architectural questions & \Cell{Track uncertainty, competing hypotheses and follow-up verification.} \\
\bottomrule
\end{tabularx}

\section{Record review and sign-off}
\Field{Prepared by, role, and revision}{Provide author, contact, version, and creation date.}
\Field{Reviewed by, date, and scope}{Record who reviewed which evidence and what remains open.}
\Field{Status and next action}{Distinguish unreviewed, draft, reviewed, approved, and superseded; attach approvals separately.}
\end{document}
